\subsection{Hệ thống nhúng}
Hệ thống nhúng là một hệ thống máy tính chuyên dụng được gắn vào thiết bị lớn hơn để thực
hiện một nhóm chức năng cố định, thường phải đáp ứng trong giới hạn thời gian, chạy liên
tục với tài nguyên hạn chế (vài trăm KB RAM, vài MB flash) và giao tiếp trực tiếp với thế
giới thực qua cảm biến và cơ cấu chấp hành. Thành phần trung tâm là \textbf{vi điều khiển}:
một vi mạch tích hợp CPU, bộ nhớ và các khối ngoại vi (GPIO, bộ định thời, PWM, SPI, ADC,
UART, Wi-Fi\ldots) trên cùng một chip. Phần mềm (firmware) được nạp vào flash và chạy ngay
khi cấp nguồn.

Một hệ thống nhúng được gọi là \textbf{thời gian thực} khi tính đúng đắn phụ thuộc cả vào
thời điểm có kết quả. Trong đề tài, thanh chắn phải đóng sau đúng 3\,s, cảnh báo cháy phải
mở thanh chắn ngay, và màn hình phải phản ánh chỗ trống trong khoảng 0,1\,s; đây là các
ràng buộc thời gian mềm (trễ một chút không gây hỏng hóc nhưng làm giảm chất lượng).

\subsection{Hệ điều hành thời gian thực FreeRTOS}
FreeRTOS là nhân hệ điều hành thời gian thực mã nguồn mở, nhỏ gọn, được Espressif tích hợp
sẵn trong ESP-IDF. Các khái niệm được sử dụng trong đề tài:
\begin{itemize}
  \item \textbf{Tác vụ (task)}: mỗi công việc là một hàm vòng lặp vô hạn có stack riêng và
        mức ưu tiên; bộ lập lịch \emph{ưu tiên, chiếm quyền} luôn chạy tác vụ sẵn sàng có
        ưu tiên cao nhất, các tác vụ cùng ưu tiên được chia thời gian theo nhịp tick.
  \item \textbf{Chặn và đánh thức}: tác vụ chờ sự kiện (gói tin đến, thông báo,
        \code{vTaskDelay}) không tốn CPU; nhờ đó CPU chỉ làm việc khi có dữ liệu.
  \item \textbf{Đồng bộ}: \emph{mutex} bảo vệ tài nguyên dùng chung trong thời gian dài
        (bus SPI của màn hình); \emph{critical section} bảo vệ vùng rất ngắn (sao chép trạng
        thái bãi); \emph{task notification} là cách nhẹ nhất để một tác vụ đánh thức tác vụ khác.
  \item \textbf{Hẹn giờ phần mềm} (\code{esp\_timer}): gọi một hàm sau một khoảng thời gian
        với độ chính xác micro giây, dùng để đóng thanh chắn.
\end{itemize}
